\renewcommand{\indexname}{索引}
% Keep only the TOC entry tied to the explicit index anchor below.
\begingroup
\makeatletter
\@mainmatterfalse
\makeatother
\begin{theindex}
\phantomsection\label{chap:index}
\addcontentsline{toc}{chapter}{索引}
\item Cartesian 积, \pageref{chap:products}, \pageref{sec:products-of-affine-varieties}
\item Cayley, Arthur, \pageref{sec:lines-on-smooth-cubic-surfaces}
\item Clebsch 三次曲面, \pageref{ex:lines-on-clebsch-cubic}
\item Clebsch, Alfred, \pageref{ex:lines-on-clebsch-cubic}
\item Cramer 法则, \pageref{lem:cramers-rule-over-ring}, \pageref{sec:dimension-of-tangent-spaces}
\item De Morgan 律, \pageref{def:topological-space}
\item Eckardt 点, \pageref{idx:eckardt-point}
\item Eisenstein 判别法, \pageref{prop:eisenstein-criterion}, \pageref{eq:quadratic-subalgebra-presentation}
\item Euclid 拓扑, \pageref{ex:containment-in-affine-plane}
\subitem 实数上的 Euclid 拓扑, \pageref{ex:euclidean-topology-on-real-line}
\item Euclid 引理
\subitem 整数的 Euclid 引理, \pageref{def:prime-element}
\subitem 多项式的 Euclid 引理, \pageref{ex:division-algorithm-fails-in-zz}
\item Euler 恒等式, \pageref{ex:projective-jacobian-criterion}
\item Grassmann 簇, \pageref{sec:grassmannians-and-plucker-maps}
\subitem 作为射影簇的 Grassmann 簇, \pageref{def:grassmannian-as-projective-variety}
\item Grassmann, Hermann, \pageref{sec:grassmannians-and-plucker-maps}
\item Gröbner 基, \pageref{ex:circle-in-plane-prime-ideal}
\item Hilbert 基定理, \pageref{thm:hilbert-basis-theorem}, \pageref{ex:projective-conic-and-affine-chart}
\item Hilbert, David, \pageref{thm:hilbert-basis-theorem}
\item Hironaka, Heisuke, \pageref{ex:rational-functions-form-algebra}, \pageref{ex:blowup-of-affine-cone}
\item Jacobian 矩阵, \pageref{def:jacobian-matrix}
\item Jacobian 判据, \pageref{def:smoothness-of-quasiprojective-variety}
\subitem 仿射超曲面的 Jacobian 判据, \pageref{ex:jacobian-criterion-for-hypersurfaces}
\subitem 仿射簇的 Jacobian 判据, \pageref{prop:jacobian-criterion-for-smoothness}
\subitem 射影超曲面的 Jacobian 判据, \pageref{ex:singularities-of-projective-hypersurfaces}
\item Krull 维数, \pageref{ex:krull-dimension-of-coordinate-ring}
\item Noether 环, \pageref{chap:irreducibility-of-affine-varieties}, \pageref{ex:fields-are-noetherian}, \pageref{def:irreducible-projective-variety}
\subitem Noether 环是分解整环, \pageref{ex:noetherian-ring-as-factorization-domain}
\subitem Noether 性不被子环保留, \pageref{ex:subring-need-not-be-noetherian}
\subitem Noether 性被取商保留, \pageref{ex:quotient-of-noetherian-ring}
\subitem 升链条件, \pageref{prop:ascending-chain-condition}
\subitem Hilbert 基定理, \pageref{thm:hilbert-basis-theorem}
\item Noether 基, \pageref{def:noether-basis}
\subitem Noether 基与维数, \pageref{cor:noether-basis-counts-dimension}
\subitem Noether 基与极小多项式, \pageref{lem:minimal-polynomials-and-noether-bases}
\subitem Noether 基与超越基, \pageref{ex:empty-transcendence-basis}
\item Noether 正规化定理, \pageref{chap:nullstellensatz-proof}, \pageref{thm:noether-normalization}, \pageref{prop:dimension-monotonicity}, \pageref{eq:transcendence-degree-difference-for-fibres}
\subitem Noether 正规化定理的证明, \pageref{thm:noether-normalization}
\item Noether, Emmy, \pageref{def:noetherian-ring}
\item Plücker 多项式, \pageref{ex:plucker-polynomial-in-six-variables}
\item Plücker 映射, \pageref{prop:plucker-maps-are-injective}
\subitem Plücker 映射是单射, \pageref{prop:plucker-maps-are-injective}
\subitem Plücker 映射的像, \pageref{prop:plucker-image-defined-by-plucker-polynomials}
\item Plücker, Julius, \pageref{def:plucker-maps}
\item Salmon, George, \pageref{sec:lines-on-smooth-cubic-surfaces}
\item Schubert 演算, \pageref{ex:lines-meeting-four-fixed-projective-lines}
\item Schubert, Hermann, \pageref{ex:lines-meeting-four-fixed-projective-lines}
\item Segre 映射, \pageref{def:segre-map}
\subitem Segre 映射的像, \pageref{prop:segre-map-injective-with-determinantal-image}
\item Segre, Corrado, \pageref{sec:products-and-segre-maps}
\item Veronese 簇, \pageref{ex:veronese-coordinate-products-depend-on-sum}
\item Veronese 映射, \pageref{def:veronese-map}
\subitem Veronese 映射是嵌入, \pageref{prop:veronese-maps-are-embeddings}
\subitem Veronese 映射的像, \pageref{prop:veronese-image-defined-by-quadratics}
\item Veronese, Giuseppe, \pageref{sec:veronese-maps}
\item Whitney 嵌入定理, \pageref{ex:tangent-space-of-single-point}
\item Zariski 拓扑
\subitem 仿射空间上的 Zariski 拓扑, \pageref{def:zariski-topology}
\subitem 射影空间上的 Zariski 拓扑, \pageref{def:zariski-topology-on-projective-subset}
\subitem 射影空间子集上的 Zariski 拓扑, \pageref{def:zariski-topology-on-projective-subset}
\item Zariski 引理, \pageref{lem:zariski}, \pageref{def:transcendence-basis}
\item Zariski, Oscar, \pageref{def:zariski-topology}
\item 爆破
\subitem 仿射空间的爆破, \pageref{def:blowup-of-affine-space-at-origin}
\subitem 仿射簇的爆破, \pageref{def:blowup-of-affine-variety-at-origin}
\item 闭包 (拓扑), \pageref{def:closure-and-dense-subset}
\item 闭嵌入
\subitem 仿射簇的闭嵌入, \pageref{ex:pullback-of-inclusion}
\item 不交并, \pageref{ex:two-dimensional-projective-space}
\item 不可约多项式, \pageref{sec:irreducible-polynomials}
\subitem 不可约多项式的刻画, \pageref{prop:irreducible-polynomial-characterization}
\item 不可约仿射簇, \pageref{sec:irreducible-affine-varieties}
\subitem 作为素理想的消没集的不可约仿射簇, \pageref{ex:irreducible-variety-in-finite-union}
\subitem 用消没理想刻画不可约仿射簇, \pageref{sec:irreducible-affine-varieties}
\subitem 不可约仿射簇的乘积, \pageref{prop:product-of-irreducible-varieties-is-irreducible}
\item 不可约分解
\subitem 不可约分解与不可约分解式, \pageref{prop:hypersurface-irreducible-components}
\subitem 仿射簇的不可约分解, \pageref{ex:quadric-intersection-complex-components}
\subitem 乘积的不可约分解, \pageref{cor:irreducible-components-of-product}
\subitem 射影簇的不可约分解, \pageref{prop:projective-irreducible-decomposition}
\subitem 拟射影簇的不可约分解, \pageref{prop:quasiprojective-irreducible-decomposition}
\item 不可约分支
\subitem 不可约分支与不同的不可约因子, \pageref{ex:quadric-intersection-real-components}
\subitem 仿射簇的不可约分支, \pageref{ex:quadric-intersection-real-components}
\subitem 乘积的不可约分支, \pageref{cor:irreducible-components-of-product}
\subitem 射影簇的不可约分支, \pageref{def:irreducible-projective-variety}
\subitem 拟射影簇的不可约分支, \pageref{ex:irreducible-components-of-dense-open-subvariety}
\item 不可约拟射影簇, \pageref{ex:closures-in-affine-line}
\item 不可约射影簇, \pageref{def:irreducible-projective-variety}
\item 不同的不可约分解式, \pageref{prop:radicals-of-principal-ideals}
\subitem 与不可约分解的关系, \pageref{prop:hypersurface-irreducible-components}
\item 不同的不可约因子, \pageref{prop:radicals-of-principal-ideals}
\subitem 与不可约分支的关系, \pageref{ex:irreducible-components-of-quadric-intersection}
\item 超曲面, \pageref{ex:dimension-of-unit-sphere}
\item 超越次数, \pageref{def:transcendence-degree}
\subitem 超越次数的界, \pageref{cor:bounds-on-transcendence-degree}
\item 超越基, \pageref{chap:dimension}, \pageref{def:transcendence-basis}
\subitem 超越基的大小相同, \pageref{prop:transcendence-bases-equal-size}
\subitem 超越基的存在性, \pageref{prop:finite-transcendence-basis-existence}
\item 超越数, \pageref{ex:pi-not-algebraic}
\item 乘法映射, \pageref{def:multiplicative-map}
\item 乘积
\subitem 仿射簇的乘积, \pageref{sec:products-of-affine-varieties}
\subsubitem 乘积与维数, \pageref{prop:dimension-additive-on-products}
\subsubitem 乘积与光滑性, \pageref{prop:tangent-space-of-product}
\subsubitem 仿射簇的乘积是仿射簇, \pageref{prop:product-of-affine-varieties-is-affine}
\subsubitem 乘积保持不可约性, \pageref{ex:reducible-cross-times-line}
\subitem 射影簇的乘积, \pageref{prop:segre-map-injective-with-determinantal-image}, \pageref{def:product-of-projective-varieties}
\subitem 拟射影簇的乘积, \pageref{def:product-of-quasiprojective-varieties}
\item 稠密子集
\subitem 仿射簇的稠密子集, \pageref{ex:pullback-of-inclusion}
\item 稠密子集 (拓扑), \pageref{def:closure-and-dense-subset}
\item 除法算法, \pageref{thm:polynomial-division}
\item 代数闭包, \pageref{sec:nullstellensatz}
\item 代数闭域, \pageref{def:algebraically-closed-field}, \pageref{lem:algebraic-extension-of-algebraically-closed-field}
\subitem 代数闭域为无限域, \pageref{thm:fundamental-theorem-of-algebra}
\item 代数基本定理, \pageref{thm:fundamental-theorem-of-algebra}
\item 代数相关性与代数独立性, \pageref{def:noether-basis}
\item 代数与几何的等价, \pageref{sec:equivalence-algebra-geometry}
\item 单项式, \pageref{def:monomial}
\subitem 单项式的次数, \pageref{def:polynomial-degree}
\item 导数, \pageref{sec:linearizations-and-tangent-spaces}
\item 地球, \pageref{fig:dimension-blue-marble}
\item 第一同构定理
\subitem $K$-代数的第一同构定理, \pageref{thm:first-isomorphism-theorem-k-algebras}
\subitem 环上代数的第一同构定理, \pageref{def:r-algebra}, \pageref{def:noether-basis}
\subitem 模的第一同构定理, \pageref{def:r-algebra}
\subitem 环的第一同构定理, \pageref{thm:first-isomorphism-theorem-for-rings}
\item 对偶
\subitem 线性映射的对偶, \pageref{ex:dual-and-double-dual-of-linear-map}
\subitem 向量空间的对偶, \pageref{def:dual-vector-space}
\subitem 向量空间基的对偶, \pageref{def:dual-of-a-basis}
\item 多项式, \pageref{def:polynomial-ring}
\subitem 次数的可加性, \pageref{prop:degree-additivity}
\subitem 多项式的次数, \pageref{def:polynomial-degree}
\subitem 多项式的齐次分量, \pageref{def:homogeneous-components}
\subitem 多项式的齐次化, \pageref{def:homogenization-of-polynomial}
\subitem 多项式的不可约性, \pageref{sec:irreducible-polynomials}
\subitem 首一多项式, \pageref{ex:square-root-two-integral}
\subitem 多项式的消没, \pageref{prop:polynomial-function-isomorphism}
\subitem 多项式与函数的区别, \pageref{prop:polynomial-function-isomorphism}
\item 多项式函数, \pageref{chap:coordinate-rings}, \pageref{ex:polynomial-functions-on-parabola}
\item 多项式函数的拉回
\subitem 拉回是一一对应, \pageref{prop:pullback-bijection}
\subitem 拉回的函子性, \pageref{cor:pullback-of-isomorphism}
\item 多项式环, \pageref{def:polynomial-ring}
\subitem 多项式环是分解整环, \pageref{def:unique-factorization-domain}
\subitem 多项式环是整环, \pageref{def:polynomial-ring}
\subitem 多项式环是 Noether 环, \pageref{cor:polynomial-rings-are-noetherian}, \pageref{sec:irreducible-decompositions}
\subitem 多项式环是唯一分解整环, \pageref{sec:unique-factorization-polynomial-rings}
\subitem 多项式环作为 $K$-代数, \pageref{ex:k-algebra-evaluation-homomorphisms}
\subitem 多项式环的递归性质, \pageref{prop:polynomial-ring-recursion}, \pageref{ex:polynomial-ring-as-module}
\item 多项式映射
\subitem 仿射簇之间的多项式映射, \pageref{chap:polynomial-maps}, \pageref{def:polynomial-map-between-affine-varieties}
\subitem 射影簇之间的多项式映射, \pageref{def:polynomial-maps-between-projective-varieties}
\subsubitem 射影簇之间的多项式映射是有限对一的, \pageref{cor:projective-polynomial-maps-are-finite-to-one}
\subsubitem 射影簇之间的多项式映射是线性的, \pageref{prop:projective-polynomial-maps-linear-up-to-isomorphism}
\item 多项式映射的拉回, \pageref{ex:polynomial-map-to-nodal-cubic}, \pageref{ex:pullback-preserves-multiplications}
\subitem 拉回与复合, \pageref{prop:pullback-preserves-compositions-and-identity}
\subitem 拉回与同构, \pageref{cor:pullback-of-isomorphism}
\subitem 拉回与恒等函数, \pageref{cor:pullback-of-isomorphism}
\subitem 拉回是 $K$-代数同态, \pageref{ex:pullback-by-polynomial-substitution}
\item 多项式组合, \pageref{def:polynomial-combinations-and-generators}
\item 二次曲面, \pageref{ex:lines-on-quadric-surfaces}
\item 泛光滑性, \pageref{ex:generic-smoothness-of-irreducible-hypersurfaces}
\item 范畴, \pageref{def:pullback-homomorphism}
\subitem 范畴等价, \pageref{sec:equivalence-algebra-geometry}
\item 仿射簇, \pageref{def:affine-variety}
\subitem 由理想定义的仿射簇, \pageref{prop:affine-varieties-defined-by-ideals}
\subitem 仿射簇的有限生成性, \pageref{cor:affine-varieties-finitely-generated}
\subitem 作为拟射影簇的仿射簇, \pageref{ex:punctured-affine-plane-is-not-affine}
\subitem 仿射簇之间的包含关系, \pageref{prop:operators-reverse-inclusions}
\subitem 仿射簇的交, \pageref{eq:coordinate-axes-union}, \pageref{def:zariski-topology}
\subitem 仿射簇的不可约性, \pageref{chap:irreducibility-of-affine-varieties}, \pageref{prop:irreducible-iff-vanishing-ideal-prime}
\subitem 仿射簇的不可约分支, \pageref{chap:irreducibility-of-affine-varieties}, \pageref{ex:quadric-intersection-real-components}
\subitem 仿射簇的不可约分解, \pageref{ex:quadric-intersection-complex-components}
\subitem 仿射簇的并, \pageref{eq:coordinate-axes-union}, \pageref{def:zariski-topology}
\item 仿射簇的切空间
\subitem 切空间与乘积, \pageref{prop:tangent-space-of-product}
\subitem 用局部环描述切空间, \pageref{prop:affine-tangent-space-dual-of-local-maximal-ideal-quotient}
\item 仿射簇的维数, \pageref{chap:dimension}
\subitem 维数与包含关系, \pageref{prop:dimension-monotonicity}
\subitem 维数与 Noether 基, \pageref{ex:noether-basis-for-circle}, \pageref{cor:noether-basis-counts-dimension}
\subitem 维数的公理, \pageref{def:axioms-for-dimension}
\subitem 由定义方程给出的维数界, \pageref{cor:dimension-lower-bound-by-number-of-equations}
\subitem 维数的界, \pageref{def:dimension-of-affine-variety}
\subitem 维数论基本定理, \pageref{chap:dimension}, \pageref{def:axioms-for-dimension}, \pageref{sec:fundamental-theorem-of-dimension-theory}, \pageref{ex:curves-with-arbitrarily-large-tangent-spaces}
\subsubitem 维数论基本定理的证明, \pageref{thm:fundamental-theorem-of-dimension-theory}
\subsubitem 维数论基本定理的加强形式, \pageref{thm:strong-fundamental-theorem-of-dimension-theory}
\subitem 维数的关键思想, \pageref{def:axioms-for-dimension}
\subitem 乘积的维数, \pageref{sec:geometric-properties-of-products}
\item 仿射开片, \pageref{def:affine-patches-and-restrictions}
\item 仿射空间, \pageref{def:affine-space}
\subitem 仿射空间的维数, \pageref{ex:dimension-of-affine-space}
\subitem 仿射空间的不可约性, \pageref{ex:affine-space-is-irreducible}
\subitem 仿射空间的光滑性, \pageref{ex:affine-space-is-smooth}
\item 仿射扭三次曲线, \pageref{ex:containment-in-affine-plane}
\item 仿射限制, \pageref{sec:affine-restrictions}, \pageref{def:affine-patches-and-restrictions}
\subitem 射影闭包的仿射限制, \pageref{prop:affine-restriction-of-projective-closure}
\item 仿射线性变换, \pageref{ex:translation-is-iso}
\item 仿射零点集, \pageref{def:vanishing-set}
\item 仿射零点理想, \pageref{def:vanishing-ideal}
\subitem 仿射消没理想是根理想, \pageref{def:radical-ideal}
\subitem 乘积的消没理想, \pageref{prop:vanishing-ideal-of-product}
\item 仿射锥, \pageref{def:affine-cone-over-projective-subset}
\item 非消没集, \pageref{prop:nonvanishing-sets-in-affine-closed-sets-are-affine}
\item 分解整环, \pageref{def:unique-factorization-domain}
\item 分式域, \pageref{def:fraction-field}
\item 根理想, \pageref{def:radical-ideal}
\subitem 根理想的刻画, \pageref{cor:characterize-radical-ideals}
\item 关联簇, \pageref{def:incidence-variety-of-lines-on-surfaces}
\item 光滑点
\subitem 仿射簇的光滑点, \pageref{chap:smoothness}, \pageref{sec:smooth-and-singular-points}
\subitem 拟射影簇的光滑点, \pageref{def:smoothness-of-quasiprojective-variety}
\item 光滑性
\subitem 仿射簇的光滑性, \pageref{ex:smoothness-is-intrinsic}
\subsubitem 光滑性与乘积, \pageref{prop:tangent-space-of-product}
\subitem 射影超曲面的光滑性, \pageref{ex:singularities-of-projective-hypersurfaces}
\subitem 拟射影簇的光滑性, \pageref{def:smoothness-of-quasiprojective-variety}
\item 函数域
\subitem 作为分式域的函数域, \pageref{def:function-field}
\subitem 作为有理函数集合的函数域, \pageref{prop:affine-function-field-equals-rational-function-field}
\subitem 不可约仿射簇的函数域, \pageref{chap:dimension}
\item 函子, \pageref{ex:cusp-not-isomorphic-to-line}
\subitem 全忠实函子, \pageref{ex:invert-pullback}
\item 环的不可约元, \pageref{def:units-irreducible-elements}
\item 环的代数扩张, \pageref{def:algebraic-and-integral-extensions}
\item 环的整扩张, \pageref{def:algebraic-and-integral-extensions}
\subitem 整扩张与有限生成模, \pageref{ex:rationals-algebraic-not-integral}
\item 环扩张中的代数元, \pageref{ex:pi-not-algebraic}
\item 环扩张中的整元, \pageref{def:algebraic-and-integral-elements}
\item 环上的代数, \pageref{chap:nullstellensatz-proof}, \pageref{def:r-algebra}
\item 环中的单位, \pageref{def:units-irreducible-elements}
\item 环中的素元, \pageref{def:prime-element}
\subitem 素元是不可约元, \pageref{def:units-irreducible-elements}
\item 基域, \pageref{sec:vanishing-set-operator}
\subitem 对基域的假设, \pageref{ex:weak-nullstellensatz}
\item 极大理想, \pageref{ex:linear-ideal-maximal}
\subitem 极大理想与点, \pageref{prop:single-points-and-maximal-ideals}
\subitem 极大理想是素理想, \pageref{cor:maximal-implies-prime}
\item 极小多项式, \pageref{ex:minimal-polynomial-of-square-root-two}
\subitem 极小多项式与 Noether 基, \pageref{lem:minimal-polynomials-and-noether-bases}
\item 既约环, \pageref{chap:coordinate-rings}, \pageref{def:nilpotent-elements-and-reduced-rings}
\subitem 既约环与模根理想的商, \pageref{prop:quotient-reduced-iff-ideal-radical}
\item 交叉相乘, \pageref{ex:cross-multiplied-projective-coordinates}
\item 交换引理, \pageref{lem:transcendence-basis-exchange}
\item 局部超曲面, \pageref{def:local-hypersurface}
\item 局部函数, \pageref{ex:local-functions-from-homogeneous-quotients}
\item 局部环, \pageref{prop:local-ring-unchanged-on-open-neighborhood}
\subitem 局部环的局部性, \pageref{prop:local-ring-unchanged-on-open-neighborhood}
\subitem 仿射簇的局部环, \pageref{lem:local-rings-of-affine-varieties-as-fractions}
\item 局部维数, \pageref{ex:affine-space-is-smooth}
\item 局部性质, \pageref{def:local-property-of-topological-spaces}
\item 矩阵子式, \pageref{cor:singular-locus-is-closed}
\item 空函数, \pageref{ex:empty-polynomial-function}
\item 蓝色弹珠照片, \pageref{fig:dimension-blue-marble}
\item 理想, \pageref{def:ideal}
\subitem 理想的陪集, \pageref{ex:cosets-of-variable-ideal}
\subitem 理想的有限生成性, \pageref{def:finitely-generated-ideal}
\subitem 理想的生成集, \pageref{def:ideal-generating-set}
\subitem 理想的积, \pageref{def:sum-and-product-of-ideals}
\subitem 理想的和, \pageref{def:sum-and-product-of-ideals}
\item 理想的根, \pageref{def:radical-of-an-ideal}
\subitem 理想的根是根理想, \pageref{prop:radicals-are-radical}
\item 连续映射, \pageref{def:continuous-map-of-topological-spaces}
\item 良序原理, \pageref{ex:complex-roots-of-unity}
\item 零点定理, \pageref{chap:varieties-and-ideals}, \pageref{prop:compose-vanishing-operators}, \pageref{ex:radical-is-an-ideal}, \pageref{thm:nullstellensatz}, \pageref{ex:weak-nullstellensatz}, \pageref{fig:real-hyperbola}, \pageref{ex:affine-space-is-irreducible}, \pageref{prop:hypersurface-irreducible-components}, \pageref{eq:quadratic-subalgebra-presentation}, \pageref{chap:nullstellensatz-proof}, \pageref{ex:radical-ideal-of-hypersurface-section}, \pageref{sec:dimension-of-tangent-spaces}, \pageref{ex:homogenizing-twisted-cubic-ideal}, \pageref{ex:projective-closure-of-affine-hyperbola}, \pageref{ex:regular-functions-on-projective-space-are-constant}, \pageref{eq:local-denominators-generate-unit-modulo-affine-ideal}
\subitem 射影零点定理, \pageref{sec:projective-nullstellensatz}
\subitem 零点定理的证明, \pageref{sec:proof-of-nullstellensatz}
\subitem 零点定理的陈述, \pageref{thm:nullstellensatz}
\subitem 零点定理名称的翻译, \pageref{sec:nullstellensatz}
\item 幂零元, \pageref{def:nilpotent-elements-and-reduced-rings}
\item 模, \pageref{chap:nullstellensatz-proof}, \pageref{def:r-module}
\subitem 模的有限生成性, \pageref{def:finitely-generated-module}
\subitem 模的生成元, \pageref{sec:module-generators}
\subitem 模的同态, \pageref{def:module-homomorphism}
\subitem 模的商, \pageref{def:r-algebra}
\subitem 由集合生成的子模, \pageref{ex:generated-submodule-versus-subalgebra}
\subitem 子模, \pageref{def:submodule}
\item 内在性质, \pageref{ex:intrinsic-ring-properties}
\item 拟射影簇, \pageref{chap:quasiprojective-varieties}, \pageref{def:quasiprojective-variety}
\subitem 拟射影簇是局部仿射的, \pageref{thm:quasiprojective-varieties-have-finite-affine-open-covers}
\item 拟射影簇的维数, \pageref{def:dimension-of-quasiprojective-variety}
\subitem 维数与稠密开子集, \pageref{ex:birational-equivalence-iff-dense-open-isomorphism}
\subitem 维数与包含关系, \pageref{def:dimension-of-quasiprojective-variety}
\subitem 维数论基本定理, \pageref{thm:fundamental-dimension-theorem-for-quasiprojective-varieties}
\item 扭三次曲线, \pageref{ex:projective-twisted-cubic}, \pageref{ex:homogenizing-twisted-cubic-ideal}, \pageref{ex:twisted-cubic-isomorphic-to-projective-line}
\item 偏导数, \pageref{def:linearization-of-polynomial}
\item 齐次多项式, \pageref{def:homogeneous-polynomial}
\item 齐次方程组的维数, \pageref{prop:dimension-of-homogeneous-systems}
\item 齐次化
\subitem 多项式的齐次化, \pageref{def:homogenization-of-polynomial}
\subitem 集合的齐次化, \pageref{def:homogenization-of-set}
\item 齐次理想, \pageref{def:homogeneous-ideal}
\item 齐次坐标环, \pageref{ex:isomorphism-does-not-preserve-homogeneous-coordinate-ring}
\item 奇点消解, \pageref{ex:rational-functions-form-algebra}
\item 奇异点
\subitem 仿射簇的奇异点, \pageref{chap:smoothness}, \pageref{sec:smooth-and-singular-points}
\subsubitem 奇异点集是闭集, \pageref{cor:singular-locus-is-closed}
\subitem 拟射影簇的奇异点, \pageref{def:smoothness-of-quasiprojective-variety}
\item 切空间
\subitem 仿射簇的切空间, \pageref{chap:smoothness}, \pageref{def:tangent-vector-and-tangent-space}
\subsubitem 切空间与梯度向量, \pageref{def:gradient-vector}
\subsubitem 切空间与同构, \pageref{cor:tangent-spaces-invariant-under-isomorphism}
\subsubitem 切空间与 Jacobian 矩阵, \pageref{def:jacobian-matrix}
\subsubitem 切空间是向量空间, \pageref{cor:tangent-space-is-vector-space}
\subsubitem 切空间维数的下界, \pageref{prop:tangent-dimension-lower-bound}
\subsubitem 切空间的内在刻画, \pageref{sec:intrinsic-tangent-spaces}, \pageref{cor:tangent-spaces-invariant-under-isomorphism}
\subitem 拟射影簇的切空间, \pageref{def:tangent-space-of-quasiprojective-variety}
\item 切向量, \pageref{chap:smoothness}, \pageref{def:tangent-vector-and-tangent-space}
\item 曲面, \pageref{ex:dimension-of-unit-sphere}
\item 曲线, \pageref{ex:dimension-of-unit-sphere}
\item 弱零点定理, \pageref{ex:weak-nullstellensatz}, \pageref{prop:weak-nullstellensatz}
\item 商
\subitem 模极大理想的商, \pageref{prop:prime-maximal-quotient-characterization}
\subitem 模素理想的商, \pageref{prop:prime-maximal-quotient-characterization}
\subitem 模根理想的商, \pageref{prop:quotient-reduced-iff-ideal-radical}
\subitem 代数的商, \pageref{def:k-subalgebra}
\subitem 模的商, \pageref{def:r-algebra}, \pageref{ex:quotient-module-is-well-defined}
\subitem 环的商, \pageref{def:quotient-ring}
\item 射影闭包, \pageref{def:projective-closure}, \pageref{ex:projective-closure-of-arbitrary-subset}
\subitem 超曲面的射影闭包, \pageref{prop:projective-closure-of-hypersurface}
\subitem 通过齐次化得到射影闭包, \pageref{prop:projective-closure-by-homogenizing-vanishing-ideal}
\item 射影簇, \pageref{chap:projective-varieties}, \pageref{ex:projective-varieties-in-projective-line}
\subitem 由理想定义的射影簇, \pageref{prop:projective-varieties-defined-by-ideals}
\subitem 由二次多项式定义的射影簇, \pageref{prop:projective-varieties-defined-by-quadratics}
\subitem 射影簇的有限生成性, \pageref{ex:classify-varieties-in-projective-line}
\subitem 射影簇的不可约性, \pageref{def:irreducible-projective-variety}
\subitem 射影簇的不可约分支, \pageref{def:irreducible-projective-variety}
\subitem 射影簇的不可约分解, \pageref{prop:projective-irreducible-decomposition}
\item 射影等价, \pageref{def:projective-equivalence}
\item 射影空间, \pageref{chap:projective-varieties}
\subitem 作为商的射影空间, \pageref{sec:projective-space}
\subitem 作为仿射空间延拓的射影空间, \pageref{ex:two-dimensional-projective-space}
\subitem 作为过原点直线集合的射影空间, \pageref{cor:projective-space-as-lines-through-origin}
\subitem 无穷远点, \pageref{prop:points-at-infinity-correspond-to-lines-through-origin}
\item 射影空间的线性子簇, \pageref{def:linear-subvariety-of-projective-space}
\item 射影零点定理, \pageref{sec:projective-nullstellensatz}
\item 射影消没, \pageref{def:projective-vanishing}
\item 射影消没集, \pageref{def:projective-vanishing-operator}
\item 射影消没理想, \pageref{sec:projective-ideal-operator}
\subitem 射影消没理想是齐次理想, \pageref{def:homogeneous-ideal}
\subitem 射影消没理想是根理想, \pageref{prop:projective-vanishing-ideal-is-homogeneous-radical}
\item 升链条件, \pageref{prop:ascending-chain-condition}
\item 首一多项式, \pageref{ex:square-root-two-integral}
\item 双对偶, \pageref{thm:vector-space-canonically-isomorphic-to-double-dual}
\item 双线性映射, \pageref{def:r-bilinear-map}
\item 双有理几何, \pageref{ex:rational-functions-form-algebra}
\item 双有理性
\subitem 不可约仿射簇的双有理性, \pageref{def:function-field}
\subitem 拟射影簇的双有理性, \pageref{prop:birational-iff-rational-function-rings-isomorphic}
\item 双直纹曲面, \pageref{ex:segre-quadric-is-doubly-ruled}
\item 素理想, \pageref{sec:prime-maximal-ideals}
\subitem 素理想与不可约仿射簇, \pageref{ex:irreducible-variety-in-finite-union}
\subitem 素理想是根理想, \pageref{prop:prime-ideals-are-radical}
\item 拓扑空间, \pageref{def:topological-space}
\item 梯度向量, \pageref{def:gradient-vector}
\item 同构
\subitem 仿射簇之间的同构, \pageref{def:isomorphism-of-affine-varieties}
\subitem 射影簇之间的同构, \pageref{def:isomorphism-of-projective-varieties}
\subitem 拟射影簇之间的同构, \pageref{def:isomorphism-of-quasiprojective-varieties}
\item 同构类, \pageref{sec:equivalence-algebra-geometry}
\item 同胚 (性质), \pageref{def:homeomorphism-of-topological-spaces}
\item 同胚 (映射), \pageref{def:homeomorphism-of-topological-spaces}
\item 椭圆曲线, \pageref{cor:irreducible-conics-isomorphic-to-projective-line}
\item 外在性质, \pageref{def:intrinsic-extrinsic-property}
\item 唯一分解整环, \pageref{def:unique-factorization-domain}
\item 维数论基本定理, \pageref{thm:fundamental-theorem-of-dimension-theory}, \pageref{thm:strong-fundamental-theorem-of-dimension-theory}, \pageref{thm:fundamental-dimension-theorem-for-quasiprojective-varieties}
\item 无关理想, \pageref{def:projective-irrelevant-ideal}
\item 无平方因子, \pageref{ex:square-free-generators-need-not-give-radical}
\item 纤维, \pageref{sec:theorem-on-fiber-dimensions}
\item 纤维维数定理
\subitem 第一部分, \pageref{thm:fiber-dimension-lower-bound}
\subitem 第二部分, \pageref{thm:fiber-dimension-equality-on-open-subset}
\item 线性化, \pageref{chap:smoothness}
\subitem 作为线性近似的线性化, \pageref{def:linearization-of-polynomial}
\subitem 一点处的线性化, \pageref{def:linearization-of-polynomial}
\subitem 多项式的线性化, \pageref{def:linearization-of-polynomial}
\subitem 乘积的线性化, \pageref{prop:tangent-space-of-product}
\item 线性投影, \pageref{ex:linear-projections-of-projective-varieties}
\item 线性映射
\subitem 射影簇之间的线性映射, \pageref{ex:linear-map-from-projective-quadric}
\subsubitem 射影簇之间的线性映射是有限对一的, \pageref{ex:normal-form-and-finite-fibres-of-projective-linear-map}
\subitem 向量空间之间的线性映射, \pageref{def:submodule}
\item 线性子空间, \pageref{def:submodule}
\item 线性组合, \pageref{def:module-linear-combination}
\item 向量空间, \pageref{chap:nullstellensatz-proof}
\item 向量空间的基, \pageref{fig:dimension-standard-basis}, \pageref{def:field-generators}
\item 向量空间的维数, \pageref{fig:dimension-standard-basis}
\item 向量空间之间的线性映射
\subitem 线性映射的行列式, \pageref{lem:determinant-of-multiplication}
\item 行列式
\subitem 行列式与 Cramer 法则, \pageref{lem:cramers-rule-over-ring}
\subitem 行列式与极小多项式, \pageref{lem:minimal-polynomials-and-noether-bases}
\item 行列式簇, \pageref{ex:segre-image-of-two-projective-lines}
\item 形式线性组合, \pageref{def:formal-linear-combination}
\item 一元多项式环, \pageref{sec:single-variable-polynomials}
\subitem 一元多项式环是主理想整环, \pageref{cor:univariate-polynomial-ring-pid}
\subitem 除法算法, \pageref{thm:polynomial-division}, \pageref{ex:dimension-of-finite-union}
\subitem 因式定理, \pageref{cor:polynomial-factor-theorem}
\subitem 有限零点定理, \pageref{cor:polynomial-finite-zeros}
\item 有理簇, \pageref{ex:fermat-cubic-not-rational}
\item 有理函数
\subitem 不可约仿射簇上的有理函数, \pageref{def:function-field}
\subitem 拟射影簇上的有理函数, \pageref{def:rational-function-on-quasiprojective-variety}
\item 有理函数的定义域, \pageref{def:ring-of-rational-functions}
\item 有理函数环, \pageref{def:ring-of-rational-functions}
\item 有理函数域, \pageref{prop:affine-function-field-equals-rational-function-field}
\item 有理映射, \pageref{def:rational-map-of-quasiprojective-varieties}
\item 有理映射的定义域, \pageref{def:rational-map-of-quasiprojective-varieties}
\item 有理映射的像, \pageref{def:rational-map-of-quasiprojective-varieties}
\item 有限生成性
\subitem 有限生成代数, \pageref{def:finitely-generated-k-algebra}
\subitem 有限生成域扩张, \pageref{cor:algebraic-field-extensions-transitive}
\subitem 有限生成理想, \pageref{def:finitely-generated-ideal}
\subitem 有限生成模, \pageref{def:finitely-generated-module}
\item 余维, \pageref{ex:dimension-of-unit-sphere}
\item 域的代数扩张
\subitem 代数扩张的传递性, \pageref{cor:algebraic-field-extensions-transitive}
\subitem 代数扩张的刻画, \pageref{cor:algebraic-field-extensions-transitive}
\item 域的生成元, \pageref{def:field-generators}
\item 域的特征, \pageref{ex:weak-nullstellensatz}
\item 域上的代数, \pageref{chap:coordinate-rings}, \pageref{ex:integers-not-k-algebra}
\subitem 作为环扩张的域上的代数, \pageref{ex:integers-not-k-algebra}
\subitem 域上代数的有限生成性, \pageref{def:finitely-generated-k-algebra}
\subitem 域上代数的第一同构定理, \pageref{thm:first-isomorphism-theorem-k-algebras}
\subitem 域上代数的生成元, \pageref{sec:generators-of-k-algebras}
\subitem 域上代数的同态, \pageref{def:k-algebra-homomorphism}
\subitem 域上代数中的多项式组合, \pageref{def:polynomial-combinations-and-generators}
\subitem 域上代数的表示, \pageref{ex:generated-subalgebra-is-smallest}
\subitem 域上代数的商, \pageref{prop:quotient-k-algebra}
\subitem 域上代数的子代数, \pageref{def:k-subalgebra}
\item 占优映射
\subitem 仿射簇之间的占优映射, \pageref{ex:pullback-of-inclusion}
\item 张量积, \pageref{chap:products}, \pageref{sec:tensor-products-of-modules}
\subitem 张量积与乘积的坐标环, \pageref{ex:coordinate-ring-of-cylinder}
\subitem 标量扩张, \pageref{ex:extension-of-scalars-has-module-structure}
\subitem 代数的张量积, \pageref{prop:tensor-product-of-algebras-is-algebra}
\subsubitem 张量积与双线性, \pageref{sec:tensor-products-and-bilinearity}
\subitem 模的张量积, \pageref{def:tensor-product-of-modules}
\subsubitem 张量积与双线性, \pageref{sec:tensor-products-and-bilinearity}
\subitem 关系, \pageref{ex:rewrite-tensor-in-three-ways}
\subitem 简单张量, \pageref{prop:tensor-product-of-algebras-is-algebra}
\item 整闭包, \pageref{ex:integral-closure-is-subring}
\item 正则函数
\subitem 仿射簇上的正则函数, \pageref{thm:regular-functions-on-affine-varieties-are-polynomial}
\subitem 拟射影簇上的正则函数, \pageref{def:regular-function-on-quasiprojective-variety}
\item 正则函数环, \pageref{def:ring-of-regular-functions}
\item 正则映射
\subitem 正则映射是连续的, \pageref{thm:regular-maps-are-zariski-continuous}
\subitem 射影簇之间的正则映射, \pageref{def:regular-maps-between-projective-varieties}
\subitem 拟射影簇之间的正则映射, \pageref{def:regular-map-of-quasiprojective-varieties}
\item 正则映射的拉回, \pageref{ex:ring-invariance-under-quasiprojective-isomorphism}
\item 正则映射的图像, \pageref{def:graph-of-quasiprojective-regular-map}
\item 直和, \pageref{fig:circle-times-line-cylinder}
\item 直线关联图, \pageref{thm:line-incidence-graph-of-smooth-cubic-is-universal}
\item 指标集, \pageref{ex:intersection-and-union-via-ideals}
\item 秩-零化度定理, \pageref{ex:invertible-affine-linear-map}, \pageref{ex:prove-dimension-function-unique}, \pageref{sec:theorem-on-fiber-dimensions}
\item 主理想, \pageref{ex:variables-ideal-not-principal}
\item 主理想的根, \pageref{prop:radicals-of-principal-ideals}
\item 主理想整环, \pageref{def:principal-ideal-domain}
\subitem 主理想整环是唯一分解整环, \pageref{def:principal-ideal-domain}
\item 子簇, \pageref{def:subvariety-of-quasiprojective-variety}
\item 子空间拓扑, \pageref{def:subspace-topology}
\item 自由模, \pageref{def:free-module-on-a-set}
\item 坐标函数, \pageref{ex:coordinate-functions}, \pageref{ex:generated-subalgebra-is-smallest}
\item 坐标环, \pageref{chap:coordinate-rings}, \pageref{def:coordinate-ring}
\subitem 坐标环与坐标函数, \pageref{cor:coordinate-ring-finitely-generated}
\subitem 作为 $K$-代数的坐标环, \pageref{prop:coordinate-ring-as-quotient-k-algebra}
\subitem 作为多项式函数环的坐标环, \pageref{ex:coordinate-ring-of-empty-set}
\subitem 作为商的坐标环, \pageref{ex:coordinate-ring-of-parabola}, \pageref{prop:coordinate-ring-as-quotient-k-algebra}
\subitem 作为正则函数环的坐标环, \pageref{ex:regular-function-ring-of-punctured-projective-conic}
\subitem 坐标环的刻画, \pageref{prop:coordinate-rings-as-finitely-generated-reduced-algebras}
\subitem 用坐标环判别同构, \pageref{cor:coordinate-rings-detect-isomorphisms}
\subitem 乘积的坐标环, \pageref{sec:coordinate-rings-of-products}
\end{theindex}
\endgroup
